\section{Task Definition and the TrajGap-Bench Protocol}
\label{sec:bench}

\subsection{Direct forecasting from an incomplete history}
\label{sec:task}

A scene at time $t$ contains $N$ pedestrians. For pedestrian $i$ the history is a pair of sequences $\{p_{i,t}\}_{t=1}^{T}$ and $\{m_{i,t}\}_{t=1}^{T}$, where $p_{i,t}\in\mathbb{R}^{2}$ is a position and $m_{i,t}\in\{0,1\}$ is a validity indicator. We use $T=8$ frames at $\Delta t=0.4$~s (3.2~s of history) and a horizon of $H=12$ frames (4.8~s), following the ETH/UCY convention. Every sample is required to contain at least one observed frame; samples that violate this are removed at load time. When $m_{i,t}=0$, the corresponding row of the position tensor holds a placeholder and must not be interpreted as a coordinate.

A direct predictor is a function
\begin{equation}
\{(\hat{y}_{i,k},\pi_{i,k})\}_{k=1}^{K}=f_{\theta}\big(\{p_{i,t},m_{i,t}\}_{i,t},\,\mathcal{C}\big),
\label{eq:task}
\end{equation}
that maps the masked history and the scene context $\mathcal{C}$ (the other agents, their masks and their geometry) to $K$ candidate futures $\hat{y}_{i,k}\in\mathbb{R}^{H\times 2}$ with mode probabilities $\pi_{i,k}\in[0,1]$, $\sum_{k}\pi_{i,k}=1$. The constraint that defines the setting is procedural rather than architectural: $f_{\theta}$ may perform any internal computation, including latent completion, but it may not call an external imputer whose output is substituted for the missing coordinates. This separates the direct setting from impute-then-predict, where the predictor receives a dense sequence and therefore strictly more information. The two settings are reported in separate tables and are never mixed.

Because raw scene coordinates differ in origin and heading, each sample is expressed in a focal-centred local frame: the origin is the last observed position of the target pedestrian, and the $x$-axis is the direction from the second-to-last observed position to the last one. Both indices are taken over \emph{observed} frames, so the heading estimate spans a gap when the history is interrupted. All agents in the scene are transformed into this frame, and the future target is expressed in it. This is the standard input normalisation for ETH/UCY, with the difference that the reference direction is defined gap-aware, which avoids deriving a heading from a placeholder row.

\subsection{Data source and the Mixed protocol}
\label{sec:mixed}

TrajGap-Bench is built on the public TrajImpute release \cite{chib2024pedestrian}, which provides, for each of five ETH/UCY scenes \cite{lerner2007crowds,pellegrini2009you} (ETH, HOTEL, UNIV, ZARA1, ZARA2), a train/validation/test split and two difficulty variants, Easy and Hard, generated from the same underlying trajectories. We do not regenerate trajectories and we do not re-split the data. For every scene and split, the benchmark forms a single Mixed set by concatenating the actor-centric samples of the two variants:
\begin{equation}
\text{Mixed} = \text{Easy} \cup \text{Hard},
\end{equation}
applied independently to train, validation and test. The union is a sample-level concatenation: no re-windowing, no re-sampling and no de-duplication. Easy and Hard labels are retained only for internal audit and are never exposed to a model; the public protocol is Mixed.

This choice has three consequences that we consider essential for a benchmark. Training sees the full range of missing-frame counts ($0$ to $7$) instead of a curated difficulty level. A single validation set determines the checkpoint for every method, so a model cannot be selected on the difficulty at which it happens to excel. And the test set preserves the official missing-level sweep, in which one latent trajectory appears under several masks. We keep that sweep because removing it would silently change the evaluation target, but it implies that test samples are not independent: two rows may describe the same underlying trajectory. Statistical comparisons must therefore be clustered by trajectory or performed with a paired bootstrap over source trajectories; we discuss this further in Section~\ref{sec:disc}.

\begin{table}[t]
\centering
\small
\caption{TrajGap-Bench Mixed splits. The test split preserves the official missing-level sweep, so its counts are sample counts rather than counts of unique source trajectories.}
\label{tab:data}
\begin{tabular}{@{}lrrr@{}}
\toprule
Scene & Train & Val & Test \\
\midrule
ETH-M   & 59{,}618 & 10{,}698 & 1{,}629 \\
HOTEL-M & 58{,}304 & 10{,}272 & 9{,}477 \\
UNIV-M  & 18{,}462 & 5{,}416  & 219{,}006 \\
ZARA1-M & 56{,}020 & 10{,}236 & 20{,}277 \\
ZARA2-M & 51{,}014 & 8{,}346  & 52{,}497 \\
\midrule
Total   & 243{,}418 & 44{,}968 & 302{,}886 \\
\bottomrule
\end{tabular}
\end{table}

Table~\ref{tab:data} lists the resulting sizes. The imbalance is inherited from the release: ETH-M contributes 0.5\% of the test samples while UNIV-M contributes 72\%. For this reason the benchmark reports both a macro average over the five scenes, which gives each scene equal weight, and a sample-weighted micro average, which reflects the distribution of the test set. Neither is a composite score, and neither replaces the per-scene table.

\subsection{Evaluation contract}
\label{sec:contract}

The contract is deliberately minimal so that it can be applied to any method.
\begin{enumerate}
\item \textbf{Training and validation.} Every method trains on Mixed-train and selects exactly one checkpoint on Mixed-val. The monitored quantity is the validation $\mathrm{minFDE}_{20}$ computed from the final prediction head; the epoch is not selected on test data and is not selected per scene.
\item \textbf{Testing.} Mixed-test is evaluated once with the selected checkpoint. No test-time tuning, thresholding or per-scene checkpoint selection is permitted.
\item \textbf{Metrics.} Two metrics define the ranking. $\mathrm{minADE}_{20}$ is the average displacement error of the candidate, among $K=20$, that is closest to the ground truth on average; $\mathrm{minFDE}_{20}$ is the final displacement error of the candidate with the smallest endpoint error. Both are in metres and lower is better. They are computed from the same set of $K=20$ candidates, the same checkpoint and the same test split, and they are never combined into a single score.
\item \textbf{Reporting.} Results are reported per scene, as a five-scene macro average, as a sample-weighted micro average and as the worst scene. Methods that do not complete all five scenes are excluded from the aggregates; partial results are not reported as benchmark conclusions.
\item \textbf{Setting separation.} Direct prediction and impute-then-predict use different inputs. A method that inserts an external imputer is evaluated in a separate setting; internal latent completion is allowed inside the direct setting and is treated as part of the model.
\end{enumerate}

$K=20$ is chosen because the release, the imputation baselines and the strongest external forecaster in our comparison all report best-of-20, and because best-of-$K$ with small $K$ saturates on crowded scenes. We additionally report top-1 metrics ($\mathrm{ADE@1}$ and $\mathrm{FDE@1}$, computed from the highest-probability mode) as diagnostics only; they are not part of the ranking, they cannot be computed for methods that output unweighted samples, and, as Section~\ref{sec:coverage} shows, they can move in the opposite direction to the primary metrics.

\subsection{Stratified diagnostics}
\label{sec:diagnostics}

A single number per scene hides the structure that the benchmark is meant to expose. We therefore report five stratifications, all computed from the validity mask and the ground truth, and all released per sample.
\begin{itemize}
\item \textbf{missing count} $m_{i}=\sum_{t}(1-m_{i,t})$, the number of absent frames in $[0,7]$. This is the coarse severity axis.
\item \textbf{terminal missingness} $\mathbb{1}[m_{i,T}=0]$, which separates gaps that end before the forecast origin from gaps that reach it.
\item \textbf{anchor lag} $T-\ell_{i}$, where $\ell_{i}$ is the index of the last observed frame. It is zero when the last frame is observed and grows as the newest trustworthy state becomes older.
\item \textbf{longest missing run}, the length of the longest contiguous gap, which distinguishes scattered dropout from a single occlusion.
\item \textbf{valid count} $T-m_{i}$, the number of observed frames, reported for completeness.
\end{itemize}

Two samples with the same missing count can differ in every other dimension: six missing frames may be one contiguous block at the end of the window, or six isolated frames scattered over the window with the last frame observed. Because velocity, heading and the trustworthiness of the last state depend on the latter structure rather than on the count, we treat the count as one axis among several. All stratifications are reported post hoc; they never define the checkpoint or the ranking.

\subsection{Scope and caveats}
The benchmark is a controlled stress test, not a measurement of real tracker statistics. The masks come from the release, which synthesises missingness; its marginal distribution is not claimed to match any particular detector or tracker. External validity must be established separately, for example with trajectories produced by a real perception stack \cite{fujii2024realtraj}. The benchmark is also restricted to third-person scenes, five ETH/UCY scenes, and a single history length. These restrictions are deliberate: they make the contract auditable, and they define the boundary within which our empirical claims hold.
